\slajdid{09_3-s008}
\begin{frame}[label=09_3-s008]
\gusttekst\setlength{\abovedisplayskip}{3pt}\setlength{\belowdisplayskip}{3pt}\setlength{\abovedisplayshortskip}{2pt}\setlength{\belowdisplayshortskip}{2pt}\centering\input{slike/tikz/s008_kapacitivnosti_dva_invertora.tex}\par\medskip\raggedright
Smatraćemo da se na ulazu u 1. Invertor nalazi idealni naponski izvor, sa beskonačnim strujnim kapacitetom, pa nam neće smetati ulazne kapacitivnosti gejta prema sorsu u oba tranzistora u PRVOM invertoru. Ali će uticati kapacitivnosti gejt prema drejnu i P i N tranzistora pošto se one pune i prazne preko tranzistora sa konačnom otpornošću. Isto tako izlaz prvog invertora menja potencijal pa dolaze do izražaja kapacitivnosti izlaza, drejna prema osnovi. Kako je izlaz 1. invertora povezan na ulaz 2. Invertora on će biti opterećen ukupnim kapacitivnostima gejtova i P i N tranzistora u 2. Invertoru.
\end{frame}
